\newif\ifglobal

%%%%%%%%%%%%%%%%%%%%%%%
%%% Compile to view %%%
%%%%%%%%%%%%%%%%%%%%%%%

%%%%%%%%%%%%%%%%%%%%%%%%%%%%%%%%%%%%%%%%%%%%%%%%%%%
%%% Readme:                                     %%%
%%% https://www.overleaf.com/read/bwdjvfjksvjy/ %%%
%%%%%%%%%%%%%%%%%%%%%%%%%%%%%%%%%%%%%%%%%%%%%%%%%%%

% >>> M?: marks a module setting number or important notice

% >>> M4: Set {./relative-path} to external files for this module <<<
\def \modulefiles {./19-DBGASSIST}
% To add an external file, use:
% (Replace '---' with actual filename)
% '\input{\modulefiles/tables/---.tex}' for tables
% '\includegraphics{\modulefiles/figures/---}' for figures
% '\subfile{\modulefiles/---}' for register files


% >>> M1: This module is in global project? <<<
% 'YES' - uncomment; 'NO' - comment out
\globaltrue

\ifglobal
    \documentclass[main\_\_CN.tex]{subfiles}
   
\else
    % Font "FandolSong-Regular" used by ctex does not have GSUB table.
    % It causes warnings that do not affect compilation and can be ignored.
    \PassOptionsToPackage{quiet}{fontspec}

    \documentclass[openany, 10pt]{book}
   
    % >>> M2: In ./00-shared/preamble-trm-module, also find
    %         settings for visibility of todo notes, watermark <<< 
    \usepackage{preamble-trm-module}


    % Enable Chinese typesetting
    \usepackage{ctex}

    % >>> M3: Temporary labels in Chinese
    %         you can add your own labels there <<<
    \myexternaldocument{./00-shared/temp-labels-cn}

\fi %\ifglobal


\setcounter{secnumdepth}{4}

% Define formatting of chapters in text
\titleformat{\chapter}
  [display]
  {\LARGE\bfseries}
  {\color{AllportsColor}第\, \thechapter\, 章}
  {1em}
  {\color{AllportsColor}}
\titlespacing*{\chapter}{0pt}{0pt}{20pt}

% Define headers
\usepackage{fancyhdr}
\usepackage{lipsum}
\renewcommand{\chaptermark}[1]{\markboth{第\, \thechapter\, 章 \  \ #1}{}}
\fancypagestyle{plain}{%
  \fancyhead{}%
  %\fancyhead[R]{\Navigation}
  \fancyhead[L]{%
    \nouppercase{%
      \ %
      \normalsize\leftmark
    }%
  }%
}
\pagestyle{plain}

% >>> M5: If this module needs any updates in
% './00-shared/preamble-trm-module.sty'
% (include additional package, etc.), contact your module PM
% or create the file './00-shared/changelog.tex' make the
% necessary updates yourself and document them in this file <<<

% Make text left-aligned and allow latex to break pages naturally, comment out for full justification
\raggedright
\raggedbottom


\begin{document}


% Sets language into which pre-defined bits of text are translated
\selectlanguage{chinese}

% Do not delete this block
\ifglobal
% Add the following front matter if
% this module is outside of global project
\else
    \InputIfFileExists{readme}{}{}
    \listoftodos
    {\let\clearpage\relax \listofdonetodos}
    \tableofcontents
    %\listoftables
    %\listoffigures
\fi


%%%%%%%%%%%%%%%%%%%%%%%%%
%%% TRM Module Begins %%%
%%%%%%%%%%%%%%%%%%%%%%%%%

\chapter{辅助调试}
\label{mod:dbgassist}
\hypertarget{dbgassist}{}

\section{概述}

\chipname{} 辅助调试模块提供一套调试功能，可在软件开发时用于调试，定位问题所在。

\section{主要特性}
\label{sec:dbgassist-features}

辅助调试模块具有以下特性：

\begin{itemize}
    \item \textbf{区域读写监测}：监测高性能处理器 (High-Performance CPU, HP CPU) 总线在限定存储器地址范围内的读写操作，若发生读写操作则触发中断。
    \item \textbf{栈指针 (SP) 监测}：监测栈指针是否超出限定范围，若超出则产生中断。
    \item \textbf{程序计数器 (PC) 记录}：记录 PC 值，以获取上一次 HP CPU 复位时的 PC 值。
    \item \textbf{总线访问记录}：记录总线访问信息，当 HP CPU 或 DMA 写入特定值时，记录此次写操作的总线类型、地址和 PC 值（仅记录 HP CPU 写操作的 PC），并将信息存储到 HP SRAM 中。
\end{itemize}

\section{功能描述}
\label{sec:dbgassist-func-descr}

\subsection{区域读写监测}

辅助调试模块监测 HP CPU 是否在某段地址范围（即区域）进行读写操作。当总线在监测地址范围内进行读写操作时，触发中断。详情请见章节 \ref{mod:riscvcpu} \textit{\nameref{mod:riscvcpu}} 中的表 \ref{tab:riscvcpu-mem-blocks}。

当 HP CPU 总线访问数据空间时，称为\textbf{数据总线}；当访问 AHB 外设空间时，称为\textbf{外设总线}。辅助调试模块可在数据总线上同时监测两个地址范围的读写操作，分别称为 HP CPU 数据总线区域 0 和 HP CPU 数据总线区域 1；也可在外设总线上监测两个地址范围，称为 HP CPU 外设总线区域 0 和 HP CPU 外设总线区域 1。区域 0 和区域 1 根据需要进行配置。

详细区域读写监测模式如下：
\begin{itemize}
    \item 数据总线读写监测
        \begin{itemize}
            \item 监测 HP CPU 数据总线在 HP CPU 数据总线区域 0 的读操作
            \item 监测 HP CPU 数据总线在 HP CPU 数据总线区域 0 的写操作
            \item 监测 HP CPU 数据总线在 HP CPU 数据总线区域 1 的读操作
            \item 监测 HP CPU 数据总线在 HP CPU 数据总线区域 1 的写操作
        \end{itemize}    
    \item 外设总线读写监测
        \begin{itemize}
            \item 监测 HP CPU 外设总线在 HP CPU 外设总线区域 0 的读操作
            \item 监测 HP CPU 外设总线在 HP CPU 外设总线区域 0 的写操作
            \item 监测 HP CPU 外设总线在 HP CPU 外设总线区域 1 的读操作
            \item 监测 HP CPU 外设总线在 HP CPU 外设总线区域 1 的写操作
        \end{itemize}
\end{itemize}

\subsection{栈指针监测}

辅助调试模块监测栈指针，以防止 HP CPU 栈溢出或错误的压栈出栈。当 HP CPU 栈指针超过监测区域的上下边界时，模块记录 PC 指针并产生中断，边界由软件配置。

详细栈指针监测模式如下：
\begin{itemize}
    \item 监测 HP CPU 栈指针是否大于监测区域上边界
    \item 监测 HP CPU 栈指针是否小于监测区域下边界
\end{itemize}

\subsection{PC 记录}

在某些情况下，软件开发者希望知道上次 HP CPU 复位时的 PC 指针。例如，当程序卡死只能复位时，开发者可能希望读取复位时的 PC 指针，以便确定程序卡死的位置并进行调试。辅助调试模块记录 HP CPU 复位时的 PC，方便开发者调试。

\subsection{CPU/DMA 总线访问记录}

辅助调试模块实时记录 HP CPU 数据总线和 DMA 总线的写操作。当总线在某个范围内发生写操作或写入特定值时，记录总线类型、地址和 PC（仅记录 HP CPU 写操作的 PC）等信息，并按照一定格式存储到 HP SRAM 中。此处提到的地址范围只能位于 HP SRAM 或外部 RAM 的范围内。

%%%%%%%%%%%%%%%%%%
%%% Interrupts %%%
%%%%%%%%%%%%%%%%%%

\section{中断}
\label{sec:dbgassist-interrupts}

\chipname{} 的辅助调试模块可以生成 BUS\_MONITOR\_INTR 中断信号并发送给\hyperref[mod:intmtrx]{中断矩阵}。这一中断信号由以下辅助调试模块的内部中断源生成：

\begin{itemize}
    \item \label{int:BUSMONITORCORE0AREADRAM00RDINT}BUS\_MONITOR\_CORE\_0\_AREA\_DRAM0\_0\_RD\_INT：当 HP CPU 数据总线在 HP CPU 数据总线区域 0 进行读操作时触发。
    \item \label{int:BUSMONITORCORE0AREADRAM00WRINT}BUS\_MONITOR\_CORE\_0\_AREA\_DRAM0\_0\_WR\_INT：当 HP CPU 数据总线在 HP CPU 数据总线区域 0 进行写操作时触发。
    \item \label{int:BUSMONITORCORE0AREADRAM01RDINT}BUS\_MONITOR\_CORE\_0\_AREA\_DRAM0\_1\_RD\_INT：当 HP CPU 数据总线在 HP CPU 数据总线区域 1 进行读操作时触发。
    \item \label{int:BUSMONITORCORE0AREADRAM01WRINT}BUS\_MONITOR\_CORE\_0\_AREA\_DRAM0\_1\_WR\_INT：当 HP CPU 数据总线在 HP CPU 数据总线区域 1 进行写操作时触发。
    \item \label{int:BUSMONITORCORE0AREAPIF0RDINT}BUS\_MONITOR\_CORE\_0\_AREA\_PIF\_0\_RD\_INT：当 HP CPU 外设总线在 HP CPU 外设总线区域 0 进行读操作时触发。
    \item \label{int:BUSMONITORCORE0AREAPIF0WRINT}BUS\_MONITOR\_CORE\_0\_AREA\_PIF\_0\_WR\_INT：当 HP CPU 外设总线在 HP CPU 外设总线区域 0 进行写操作时触发。
    \item \label{int:BUSMONITORCORE0AREAPIF1RDINT}BUS\_MONITOR\_CORE\_0\_AREA\_PIF\_1\_RD\_INT：当 HP CPU 外设总线在 HP CPU 外设总线区域 1 进行读操作时触发。
    \item \label{int:BUSMONITORCORE0AREAPIF1WRINT}BUS\_MONITOR\_CORE\_0\_AREA\_PIF\_1\_WR\_INT：当 HP CPU 外设总线在 HP CPU 外设总线区域 1 进行写操作时触发。
    \item \label{int:BUSMONITORCORE0SPSPILLMININT}BUS\_MONITOR\_CORE\_0\_SP\_SPILL\_MIN\_INT：当 HP CPU 栈指针小于 HP CPU 栈指针监测区域的下边界时触发。
    \item \label{int:BUSMONITORCORE0SPSPILLMAXINT}BUS\_MONITOR\_CORE\_0\_SP\_SPILL\_MAX\_INT：当 HP CPU 栈指针大于 HP CPU 栈指针监测区域的上边界时触发。
\end{itemize}

\begin{tiplisting}
对于中断、中断信号、中断源 及其关系的详细描述，可以查看章节 \ref{mod:intmtrx} \textit{\nameref{mod:intmtrx}} > Section \ref{sec:intmtx-int-terminology} \textit{\nameref{sec:intmtx-int-terminology}}。
\end{tiplisting}

每个中断源都有一组通用的配置寄存器，详见章节 \hyperref[interrupt-config-registers]{\textit{中断配置寄存器}}。具体寄存器名称请查看 \ref{sec:dbgassist-reg-summ} \textit{\nameref{sec:dbgassist-reg-summ}}。


\section{配置流程}
\label{sec:dbgassist-prog-procedures}

\subsection{区域监测和栈指针监测配置}

区域监测和栈指针监测的配置流程如下。

\begin{enumerate}
    \item 配置监测区域和栈指针。
        \begin{itemize}
            \item HP CPU 数据总线区域 0 由 \hyperref[regdesc:BUSMONITORCORE0AREADRAM00MINREG]{BUS\_MONITOR\_CORE\_0\_AREA\_DRAM0\_0\_MIN\_REG} 和 \hyperref[regdesc:BUSMONITORCORE0AREADRAM00MAXREG]{BUS\_MONITOR\_CORE\_0\_AREA\_DRAM0\_0\_MAX\_REG} 配置。
            \item HP CPU 数据总线区域 1 由 \hyperref[regdesc:BUSMONITORCORE0AREADRAM01MINREG]{BUS\_MONITOR\_CORE\_0\_AREA\_DRAM0\_1\_MIN\_REG} 和 \hyperref[regdesc:BUSMONITORCORE0AREADRAM01MAXREG]{BUS\_MONITOR\_CORE\_0\_AREA\_DRAM0\_1\_MAX\_REG} 配置。
            \item HP CPU 外设总线区域 0 由 \hyperref[regdesc:BUSMONITORCORE0AREAPIF0MINREG]{BUS\_MONITOR\_CORE\_0\_AREA\_PIF\_0\_MIN\_REG} 和 \hyperref[regdesc:BUSMONITORCORE0AREAPIF0MAXREG]{BUS\_MONITOR\_CORE\_0\_AREA\_PIF\_0\_MAX\_REG} 配置。
            \item HP CPU 外设总线区域 1 由 \hyperref[regdesc:BUSMONITORCORE0AREAPIF1MINREG]{BUS\_MONITOR\_CORE\_0\_AREA\_PIF\_1\_MIN\_REG} 和 \hyperref[regdesc:BUSMONITORCORE0AREAPIF1MAXREG]{BUS\_MONITOR\_CORE\_0\_AREA\_PIF\_1\_MAX\_REG} 配置。
            \item HP CPU 栈指针监测区域由 \hyperref[regdesc:BUSMONITORCORE0SPMINREG]{BUS\_MONITOR\_CORE\_0\_SP\_MIN\_REG} 和 \hyperref[regdesc:BUSMONITORCORE0SPMAXREG]{BUS\_MONITOR\_CORE\_0\_SP\_MAX\_REG} 配置。
        \end{itemize}
    \item 配置 \hyperref[regdesc:BUSMONITORCORE0INTRENAREG]{BUS\_MONITOR\_CORE\_0\_INTR\_ENA\_REG} 使能不同模式的中断。
    \item 配置 \hyperref[regdesc:BUSMONITORCORE0MONTRENAREG]{BUS\_MONITOR\_CORE\_0\_MONTR\_ENA\_REG} 使能不同的监测模式，可同时使能。
    \item 读取 \hyperref[regdesc:BUSMONITORCORE0INTRRAWREG]{BUS\_MONITOR\_CORE\_0\_INTR\_RAW\_REG} 查询不同模式的原始中断状态。
    \item 配置 \hyperref[regdesc:BUSMONITORCORE0INTRCLRREG]{BUS\_MONITOR\_CORE\_0\_INTR\_CLR\_REG} 清除不同模式的中断。
\end{enumerate}

若要监测 HP CPU 数据总线地址 A 到地址 B 范围内的写操作，可通过 HP CPU 数据总线区域 0 或区域 1 进行监测。以 HP CPU 数据总线区域 0 为例：

\begin{enumerate}
    \item 配置 \hyperref[fielddesc:BUSMONITORCORE0RCDPDEBUGEN]{BUS\_MONITOR\_CORE\_0\_RCD\_PDEBUGEN} 为 1，使能 HP CPU 更新 PC 信号给辅助调试模块。
    \item 配置 \hyperref[regdesc:BUSMONITORCORE0AREADRAM00MINREG]{BUS\_MONITOR\_CORE\_0\_AREA\_DRAM0\_0\_MIN\_REG} 为 A 地址。
    \item 配置 \hyperref[regdesc:BUSMONITORCORE0AREADRAM00MAXREG]{BUS\_MONITOR\_CORE\_0\_AREA\_DRAM0\_0\_MAX\_REG} 为 B 地址。
    \item 配置 \hyperref[regdesc:BUSMONITORCORE0INTRENAREG]{BUS\_MONITOR\_CORE\_0\_INTR\_ENA\_REG} bit[1] 使能 HP CPU 数据总线区域 0 写监测中断。
    \item 配置 \hyperref[regdesc:BUSMONITORCORE0MONTRENAREG]{BUS\_MONITOR\_CORE\_0\_MONTR\_ENA\_REG} bit[1] 使能 HP CPU 数据总线区域 0 写监测。
    \item 配置中断矩阵将 BUS\_MONITOR\_INTR 映射到 HP CPU 中断（参考章节 \ref{mod:intmtrx} \textit{\nameref{mod:intmtrx}}）。
    \item 触发中断后：
    \begin{itemize}
        \item 读取 \hyperref[regdesc:BUSMONITORCORE0INTRRAWREG]{BUS\_MONITOR\_CORE\_0\_INTR\_RAW\_REG} 确认中断源。
        \item 如果中断源是 HP CPU 区域监测，读取 \hyperref[fielddesc:BUSMONITORCORE0AREAPC]{BUS\_MONITOR\_CORE\_0\_AREA\_PC} 获取触发中断时的 HP CPU PC 值，读取 \hyperref[fielddesc:BUSMONITORCORE0AREASP]{BUS\_MONITOR\_CORE\_0\_AREA\_SP} 获取触发中断时的 HP CPU 栈指针。
        \item 如果中断源是栈指针监测，读取 \hyperref[fielddesc:BUSMONITORCORE0SPPC]{BUS\_MONITOR\_CORE\_0\_SP\_PC} 获取触发中断时的 HP CPU PC 值。
        \item 配置 \hyperref[regdesc:BUSMONITORCORE0INTRCLRREG]{BUS\_MONITOR\_CORE\_0\_INTR\_CLR\_REG} 对应位为 1 清除中断。
    \end{itemize}
\end{enumerate}

\subsection{PC 记录配置}

将 \hyperref[fielddesc:BUSMONITORCORE0RCDPDEBUGEN]{BUS\_MONITOR\_CORE\_0\_RCD\_PDEBUGEN} 配置为 1，以使能 HP CPU 更新 PC 信号给辅助调试模块。如果 \hyperref[fielddesc:BUSMONITORCORE0RCDRECORDEN]{BUS\_MONITOR\_CORE\_0\_RCD\_RECORDEN} 也配置为 1，\hyperref[regdesc:BUSMONITORCORE0RCDPDEBUGPCREG]{BUS\_MONITOR\_CORE\_0\_RCD\_PDEBUGPC\_REG} 将记录 HP CPU 的 PC 信号，同时 \hyperref[regdesc:BUSMONITORCORE0RCDPDEBUGSPREG]{BUS\_MONITOR\_CORE\_0\_RCD\_PDEBUGSP\_REG} 将记录 HP CPU 的栈指针，否则保留原值。

当 HP CPU 复位时，\hyperref[regdesc:BUSMONITORCORE0RCDENREG]{BUS\_MONITOR\_CORE\_0\_RCD\_EN\_REG} 将被复位，但 \hyperref[regdesc:BUSMONITORCORE0RCDPDEBUGPCREG]{BUS\_MONITOR\_CORE\_0\_RCD\_PDEBUGPC\_REG} 和 \hyperref[regdesc:BUSMONITORCORE0RCDPDEBUGSPREG]{BUS\_MONITOR\_CORE\_0\_RCD\_PDEBUGSP\_REG} 不会被复位，因此这两个寄存器将保留 HP CPU 复位时的 PC 值和栈指针。

\subsection{CPU/DMA 总线访问记录配置}

\subsubsection{总线访问记录 1 配置}
\label{sec:hp-cpu-access-mon}

总线访问记录 1 包括 HP CPU 总线访问记录（HP SRAM 和外部 RAM 监控功能）和 DMA 总线访问记录（HP SRAM 监控功能）。其配置流程如下：

\begin{enumerate}
    \item 配置监测地址范围：配置 \hyperref[regdesc:MEMMONITORLOGMINREG]{MEM\_MONITOR\_LOG\_MIN\_REG} 和 \hyperref[regdesc:MEMMONITORLOGMAXREG]{MEM\_MONITOR\_LOG\_MAX\_REG} 确定地址范围。对于 HP CPU，监测地址范围需在 HP SRAM (0x4080\_0000 \verb+~+ 0x4084\_FFFF) 或外部 RAM (0x4200\_0000 \verb+~+ 0x43FF\_FFFF) 的地址区间内；对于 DMA，监测地址范围需在 HP SRAM (0x4080\_0000 \verb+~+ 0x4084\_FFFF) 的地址区间内。
    \item 配置监测模式 (\hyperref[fielddesc:MEMMONITORLOGMODE]{MEM\_MONITOR\_LOG\_MODE})：
        \begin{itemize}
            \item 写监测（监测总线写操作）
            \item 字 (word) 监测（监测总线写特定 word）
            \item 半字 (halfword) 监测（监测总线写特定 halfword）
            \item 字节 (byte) 监测（监测总线写特定 byte）
        \end{itemize}
    \item 配置需要监测的特殊值：
        \begin{itemize}
            \item 在 word 监测模式下，\hyperref[regdesc:MEMMONITORLOGCHECKDATAREG]{MEM\_MONITOR\_LOG\_CHECK\_DATA\_REG} 为监测的特定 word。
            \item 在 halfword 监测模式下，\hyperref[regdesc:MEMMONITORLOGCHECKDATAREG]{MEM\_MONITOR\_LOG\_CHECK\_DATA\_REG}[15:0] 为监测的特定 halfword。
            \item 在 byte 监测模式下，\hyperref[regdesc:MEMMONITORLOGCHECKDATAREG]{MEM\_MONITOR\_LOG\_CHECK\_DATA\_REG}[7:0] 为监测的特定 byte。
            \item \hyperref[regdesc:MEMMONITORLOGDATAMASKREG]{MEM\_MONITOR\_LOG\_DATA\_MASK\_REG} 用于屏蔽 \hyperref[regdesc:MEMMONITORLOGCHECKDATAREG]{MEM\_MONITOR\_LOG\_CHECK\_DATA\_REG} 的相应 byte。当某一 byte 被屏蔽，表示该 byte 可为任意值。例如，在 word 监测模式下，若 \hyperref[regdesc:MEMMONITORLOGCHECKDATAREG]{MEM\_MONITOR\_LOG\_CHECK\_DATA\_REG} 配置为 0x01020304，\hyperref[regdesc:MEMMONITORLOGDATAMASKREG]{MEM\_MONITOR\_LOG\_DATA\_MASK\_REG} 配置为 0x1，此时只要总线写 0x010203XX，都会被记录。
        \end{itemize}
    \item 配置记录信息的存储地址范围：
        \begin{itemize}
            \item \hyperref[regdesc:MEMMONITORLOGMEMSTARTREG]{MEM\_MONITOR\_LOG\_MEM\_START\_REG} 和 \hyperref[regdesc:MEMMONITORLOGMEMENDREG]{MEM\_MONITOR\_LOG\_MEM\_END\_REG} 为存储地址范围配置寄存器。记录信息的存储地址范围为 0x4080\_0000 \verb+~+ 0x4084\_FFFF。
            \item 置位 \hyperref[regdesc:MEMMONITORLOGMEMADDRUPDATEREG]{MEM\_MONITOR\_LOG\_MEM\_ADDR\_UPDATE\_REG}，使 \hyperref[regdesc:MEMMONITORLOGMEMCURRENTADDRREG]{MEM\_MONITOR\_LOG\_MEM\_CURRENT\_ADDR\_REG} 更新为 \hyperref[regdesc:MEMMONITORLOGMEMSTARTREG]{MEM\_MONITOR\_LOG\_MEM\_START\_REG}。
            \item 配置辅助调试模块对 HP SRAM 的使用权限。只有开启辅助调试模块对 HP SRAM 的使用权限，才能访问 HP SRAM。关于如何配置，详情见章节 \ref{mod:permctrl} \textit{\nameref{mod:permctrl}}。
        \end{itemize}
    \item \label{other:dbgassist-item-log-config-step5} 配置记录信息的存储模式，分为 loop 模式和非 loop 模式。
        \begin{itemize}
            \item 在 loop 模式下，记录信息循环写入配置地址范围。当写到结束地址时，回到起始地址开始写，覆盖之前记录的信息。置位 \hyperref[fielddesc:MEMMONITORLOGMEMLOOPENABLE]{MEM\_MONITOR\_LOG\_MEM\_LOOP\_ENABLE} 使能 loop 模式。\\
            例如，存储记录信息的地址范围为 0 \verb+~+ 4 ，在总线访问过程中进行 10 次写操作。第 5 次写操作写到地址 4 后，第 6 次写操作会回到地址 0 开始写，第 6 \verb+~+ 10 次写操作会覆盖之前的第 1 \verb+~+ 5 次写操作的数据。
            \item 在非 loop 模式下，当记录信息写到结束地址后，会停留在结束地址，不再写入记录信息，不会覆盖最初的记录信息，并丢弃后续的记录信息。清零 \hyperref[fielddesc:MEMMONITORLOGMEMLOOPENABLE]{MEM\_MONITOR\_LOG\_MEM\_LOOP\_ENABLE} 使能非 loop 模式。\\
            例如，存储记录信息的地址范围为 0 \verb+~+ 4 ，在总线访问过程中进行 10 次写操作。第 5 次写操作写到地址 4 后，第 6 \verb+~+ 10 次写操作的地址会停留在地址 4，但不再进行写操作，因此地址范围 0 \verb+~+ 4 保留前 5 次写操作的信息，并丢弃第 6 \verb+~+ 10 次写操作的信息。
        \end{itemize}
    \item 配置总线使能
    \begin{itemize}
        \item 配置 \hyperref[fielddesc:MEMMONITORLOGCOREENA]{MEM\_MONITOR\_LOG\_CORE\_ENA} 和 \hyperref[fielddesc:MEMMONITORLOGDMA0ENA]{MEM\_MONITOR\_LOG\_DMA\_0\_ENA} 使能 HP CPU 和 DMA 总线访问记录，可同时使能。
    \end{itemize}
\end{enumerate}
    
辅助调试模块将记录信息缓存到内部 buffer A 中，然后从 buffer 中取出记录信息存储到配置的存储区域。当监测行为连续触发时，产生的记录数据包可能会填满 buffer，无法缓存新的记录数据包，此时模块会舍弃无法及时缓存到 buffer 的记录数据包，并在 buffer 不满时缓存一个 LOST 数据包作为替代。此时仅能获取 LOST 数据包生成时所舍弃的记录数据包的总线类型，但无法得知之前舍弃了多少个记录数据包。

当总线访问记录结束后，需要从内存中读取记录数据并进行解码。记录数据包含三种包格式：HP CPU 数据包、DMA 数据包和 LOST 数据包。前两种数据包分别对应 HP CPU 数据总线记录信息和 DMA 数据总线记录信息，三种包格式如表 \ref{tab:dbgassist-hpcpu-package}、\ref{tab:dbgassist-dma-package} 和 \ref{tab:dbgassist-lost-package} 所示：
    
    \begin{longtable}[c]{|L{3cm}|L{3cm}|L{3cm}|L{3cm}|L{3cm}|}
    \caption{HP CPU 包格式}\label{tab:dbgassist-hpcpu-package}
    \\\hline
    \rowcolor{lightgray}
    \textbf{Bit[63:33]} & \textbf{Bit[32]} & \textbf{Bit[31:3]} & \textbf{Bit[2:1]} & \textbf{Bit[0]} \\\hline
    pc\_offset & anchored(1) & addr\_offset & format & anchored(0) \\\hline
    \end{longtable}

    \begin{longtable}[c]{|L{2cm}|L{2cm}|L{2cm}|L{2cm}|}
    \caption{DMA 包格式}\label{tab:dbgassist-dma-package}
    \\\hline
    \rowcolor{lightgray}
    \textbf{Bit[31:8]} & \textbf{Bit[7:3]} & \textbf{Bit[2:1]} & \textbf{Bit[0]} \\\hline
    addr\_offset & dma\_source & format & anchored(0) \\\hline
    \end{longtable}
    
    \begin{longtable}[c]{|L{2cm}|L{2cm}|L{2cm}|L{2cm}|}
    \caption{LOST 包格式}\label{tab:dbgassist-lost-package}
    \\\hline
    \rowcolor{lightgray}
    \textbf{Bit[31:7]} & \textbf{Bit[6:3]} & \textbf{Bit[2:1]} & \textbf{Bit[0]} \\\hline
    保留 & lost\_source & format & anchored(0) \\\hline
    \end{longtable}
    
    从包格式可以看出，HP CPU 数据包共 64 位，其他数据包均为 32 位。以下为各个域的含义：
    
    \begin{itemize}
    \item \textbf{format} 表示数据包的类型：
\begin{itemize}
    \item 0：HP CPU 数据包
    \item 1：DMA 数据包
    \item 2：保留
    \item 3：LOST 数据包
\end{itemize}
\item \textbf{pc\_offset} 记录 HP CPU 的 PC 指针偏移量，实际 PC = pc\_offset + 0x0000\_0000。
\item \textbf{addr\_offset} 记录此次写操作的地址偏移，实际地址 = addr\_offset + \{\hyperref[regdesc:MEMMONITORLOGMINREG]{MEM\_MONITOR\_LOG\_MIN\_REG}[31:2], 2'b0\}。
\item \textbf{dma\_source} 记录 DMA 的访问来源，详见表 \ref{tab:assistdbg-dma-source}，值 16 \verb+~+ 31 对应的来源参考 \ref{mod:dma} \textit{\nameref{mod:dma}}。
\item \textbf{lost\_source} 记录 LOST 数据包生成时舍弃的数据包：
\begin{itemize}
    \item Bit[4:3]：
    \begin{itemize}
        \item 0：未舍弃 HP CPU 数据包
        \item 3：舍弃 HP CPU 数据包
        \item 其他值：保留
    \end{itemize}
    \item Bit[5]：
    \begin{itemize}
        \item 0：未舍弃 DMA 数据包
        \item 1：舍弃 DMA 数据包
    \end{itemize}
    \item Bit[6]：保留
\end{itemize}
\item \textbf{anchored} 表示该 32 位数据在数据包中的位置：
\begin{itemize}
    \item 0：该 32 位数据是数据包的低 32 位
    \item 1：该 32 位数据是数据包的高 32 位
\end{itemize}
\end{itemize}

\begin{longtable}{|l| m{8cm} |}
    \caption{DMA 访问来源}\label{tab:assistdbg-dma-source}
    \\\hline
    \rowcolor{lightgray}
    \textbf{值} & \textbf{来源} \\\hline
    \endfirsthead \hline \rowcolor{lightgray}
    \textbf{值} & \textbf{来源} \\\hline
    \endhead \endfoot \hline \endlastfoot
    0 & HP CPU \\\hline
    1 & 保留 \\\hline
    2 & 保留 \\\hline
    3 & 保留 \\\hline
    4 & 保留 \\\hline
    5 & MEM\_MONITOR \\\hline
    6 & TRACE \\\hline 
    7 & 保留 \\\hline
    8 & PSRAM\_MEM\_MONITOR \\\hline
    9 \verb+~+ 15 & 保留 \\\hline
    16 \verb+~+ 31 & 详见章节 \ref{mod:dma} \textit{\nameref{mod:dma}} > 表 \ref{tab:gdma-ahb-peri-sel} 中 0 \verb+~+ 15 对应的外设。例如，16 对应的来源为该表中 0 对应的外设，17 对应的来源为该表中 1 对应的外设，以此类推 \\\hline
\end{longtable}

模块内部 buffer 的数据宽度是 32 位。若同时使能 HP CPU 和 DMA 总线访问记录，并且在某一时刻同时生成记录数据，则先将 HP CPU 数据包缓存到 buffer 中，再缓存 DMA 数据包。辅助调试模块会自动将缓存数据从 buffer 取出，并将取出的数据以 32 位数据宽度存入指定范围内存中。

当启用 loop 模式时，若在配置的存储地址范围内循环若干次，可能存在残留数据干扰解析，即 HP CPU 数据包的低 32 位数据被覆盖，使得高 32 位数据成为残留数据。因此，确定第一个有效数据包的位置时，必须过滤可能存在的残留数据。读取 \hyperref[regdesc:MEMMONITORLOGMEMCURRENTADDRREG]{MEM\_MONITOR\_LOG\_MEM\_CURRENT\_ADDR\_REG} 可确定数据包的起始位置，检查数据包 anchored 位的数值，0 则保留，1 则舍弃。

数据包解析流程如下：

\begin{itemize}
    \item 读取 \hyperref[fielddesc:MEMMONITORLOGMEMFULLFLAG]{MEM\_MONITOR\_LOG\_MEM\_FULL\_FLAG} 确定数据是否溢出配置范围。
    \begin{itemize}
        \item 如果未溢出，则数据读取的地址范围为 \hyperref[regdesc:MEMMONITORLOGMEMSTARTREG]{MEM\_MONITOR\_LOG\_MEM\_START\_REG} \verb+~+ \hyperref[regdesc:MEMMONITORLOGMEMCURRENTADDRREG]{MEM\_MONITOR\_LOG\_MEM\_CURRENT\_ADDR\_REG} $-$ 4；
        \item 若溢出且开启了 loop 模式，则数据读取的地址范围为 \hyperref[regdesc:MEMMONITORLOGMEMCURRENTADDRREG]{MEM\_MONITOR\_LOG\_MEM\_CURRENT\_ADDR\_REG} \verb+~+ \hyperref[regdesc:MEMMONITORLOGMEMENDREG]{MEM\_MONITOR\_LOG\_MEM\_END\_REG} 和 \hyperref[regdesc:MEMMONITORLOGMEMSTARTREG]{MEM\_MONITOR\_LOG\_MEM\_START\_REG} \verb+~+ \hyperref[regdesc:MEMMONITORLOGMEMCURRENTADDRREG]{MEM\_MONITOR\_LOG\_MEM\_CURRENT\_ADDR\_REG} $-$ 4；
        \item 若溢出且未开启 loop 模式，则数据读取的地址范围为 \hyperref[regdesc:MEMMONITORLOGMEMSTARTREG]{MEM\_MONITOR\_LOG\_MEM\_START\_REG} \verb+~+ \hyperref[regdesc:MEMMONITORLOGMEMENDREG]{MEM\_MONITOR\_LOG\_MEM\_END\_REG}。
    \end{itemize}
    \item 从起始地址处开始读取数据并解析，每次读取 32 位。
\end{itemize}

数据解析完成之后需要通过置位 \hyperref[fielddesc:MEMMONITORCLRLOGMEMFULLFLAG]{MEM\_MONITOR\_CLR\_LOG\_MEM\_FULL\_FLAG} 清除 \hyperref[fielddesc:MEMMONITORLOGMEMFULLFLAG]{MEM\_MONITOR\_LOG\_MEM\_FULL\_FLAG} 标志位。

\subsubsection{总线访问记录 2 配置}

总线访问记录 2 包括 DMA 总线访问记录（外部 RAM 监控功能）。其配置流程如下：

\begin{enumerate}
    \item 配置监测地址范围：配置 \hyperref[regdesc:MEMMONITORLOGMINREG]{MEM\_MONITOR\_LOG\_MIN\_REG} 和 \hyperref[regdesc:MEMMONITORLOGMAXREG]{MEM\_MONITOR\_LOG\_MAX\_REG} 确定地址范围。监测地址范围需在 0x4200\_0000 \verb+~+ 0x43FF\_FFFF 之间。
    \item 配置 \hyperref[fielddesc:MEMMONITORLOGMODE]{MEM\_MONITOR\_LOG\_MODE} 选择监测模式：
        \begin{itemize}
            \item 写监测（监测总线写操作）
            \item 字 (word) 监测（监测总线写特定 word）
            \item 半字 (halfword) 监测（监测总线写特定 halfword）
            \item 字节 (byte) 监测（监测总线写特定 byte）
        \end{itemize}
    \item 配置需要监测的特殊值：
        \begin{itemize}
            \item 在 word 监测模式下，\hyperref[regdesc:MEMMONITORLOGCHECKDATAREG]{MEM\_MONITOR\_LOG\_CHECK\_DATA\_REG} 为监测的特定 word。
            \item 在 halfword 监测模式下，\hyperref[regdesc:MEMMONITORLOGCHECKDATAREG]{MEM\_MONITOR\_LOG\_CHECK\_DATA\_REG}[15:0] 为监测的特定 halfword。
            \item 在 byte 监测模式下，\hyperref[regdesc:MEMMONITORLOGCHECKDATAREG]{MEM\_MONITOR\_LOG\_CHECK\_DATA\_REG}[7:0] 为监测的特定 byte。
            \item \hyperref[regdesc:MEMMONITORLOGDATAMASKREG]{MEM\_MONITOR\_LOG\_DATA\_MASK\_REG} 用于屏蔽 \hyperref[regdesc:MEMMONITORLOGCHECKDATAREG]{MEM\_MONITOR\_LOG\_CHECK\_DATA\_REG} 的相应 byte。当某一 byte 被屏蔽时，该 byte 可为任意值。例如，在 word 监测中，若 \hyperref[regdesc:MEMMONITORLOGCHECKDATAREG]{MEM\_MONITOR\_LOG\_CHECK\_DATA\_REG} 配置为 0x01020304，\hyperref[regdesc:MEMMONITORLOGDATAMASKREG]{MEM\_MONITOR\_LOG\_DATA\_MASK\_REG} 配置为 0x1，则只要总线写 0x010203XX 都会被记录。
        \end{itemize}
    \item 配置记录信息的存储地址范围：
        \begin{itemize}
            \item \hyperref[regdesc:MEMMONITORLOGMEMSTARTREG]{MEM\_MONITOR\_LOG\_MEM\_START\_REG} 和 \hyperref[regdesc:MEMMONITORLOGMEMENDREG]{MEM\_MONITOR\_LOG\_MEM\_END\_REG} 为存储地址范围配置寄存器。记录信息的存储地址范围为 0x4080\_0000 \verb+~+ 0x4084\_FFFF。
            \item 置位 \hyperref[regdesc:MEMMONITORLOGMEMADDRUPDATEREG]{MEM\_MONITOR\_LOG\_MEM\_ADDR\_UPDATE\_REG}，使 \hyperref[regdesc:MEMMONITORLOGMEMCURRENTADDRREG]{MEM\_MONITOR\_LOG\_MEM\_CURRENT\_ADDR\_REG} 更新为 \hyperref[regdesc:MEMMONITORLOGMEMSTARTREG]{MEM\_MONITOR\_LOG\_MEM\_START\_REG}。
            \item 配置辅助调试模块对 HP SRAM 的使用权限。只有开启辅助调试模块对 HP SRAM 的使用权限才能访问 HP SRAM。关于如何配置，详情见章节 \ref{mod:permctrl} \textit{\nameref{mod:permctrl}}。
        \end{itemize}
    \item 配置记录信息的存储模式，分为 loop 模式和非 loop 模式。
        \begin{itemize}
            \item 在 loop 模式下，记录信息循环写入配置地址范围。当写到结束地址时，回到起始地址开始写，覆盖之前记录的信息。置位 \hyperref[fielddesc:MEMMONITORLOGMEMLOOPENABLE]{MEM\_MONITOR\_LOG\_MEM\_LOOP\_ENABLE} 启用 loop 模式。
            \item 在非 loop 模式下，当记录信息写到结束地址后，停留在结束地址，不再写入记录信息，不会覆盖最开始的记录信息，并丢弃后续记录信息。清零 \hyperref[fielddesc:MEMMONITORLOGMEMLOOPENABLE]{MEM\_MONITOR\_LOG\_MEM\_LOOP\_ENABLE} 启用非 loop 模式。\\
            \item 示例参考小节 \ref{sec:hp-cpu-access-mon} 的步骤 \ref{other:dbgassist-item-log-config-step5}。
        \end{itemize}
\item 配置 \hyperref[fielddesc:MEMMONITORLOGDMA0ENA]{MEM\_MONITOR\_LOG\_DMA\_0\_ENA} 使能 DMA\_0 通道组总线访问记录。
\end{enumerate}

辅助调试模块将 DMA 记录信息缓存到内部 buffer B，然后从 buffer 中取出记录信息并存储到配置的存储区域。当监测行为连续触发时，产生的记录数据包可能使 buffer 满，无法缓存新的记录数据包，此时模块会舍弃无法及时缓存的记录数据包，并在 buffer 不满时缓存一个 DMA LOST 数据包作为替代。此时仅能知道 DMA LOST 数据包生成时所舍弃的记录数据包的总线类型，但无法得知之前舍弃了多少个记录数据包。

当总线访问记录结束后，需要从内存中读取记录数据并进行解码。记录数据仅有两种包格式：DMA\_0 数据包和 DMA LOST 数据包。DMA\_0 数据包对应 DMA\_0 通道组数据总线记录信息，两种包格式如表 \ref{tab:dbgassist-dma0-package} 和 \ref{tab:dbgassist-dma-lost-package} 所示：

\begin{longtable}[c]{|L{2cm}|L{2cm}|L{2cm}|L{2cm}|}
    \caption{DMA\_0 包格式}\label{tab:dbgassist-dma0-package}
    \\\hline
    \rowcolor{lightgray}
    \textbf{Bit[31:7]} & \textbf{Bit[6:2]} & \textbf{Bit[1]} & \textbf{Bit[0]} \\\hline
    addr\_offset & dma\_source & format & anchored(0)  \\\hline
\end{longtable}

\begin{longtable}[c]{|L{2cm}|L{2cm}|L{2cm}|L{2cm}|}
    \caption{DMA LOST 包格式}\label{tab:dbgassist-dma-lost-package}
    \\\hline
    \rowcolor{lightgray}
    \textbf{Bit[31:3]} & \textbf{Bit[2]} & \textbf{Bit[1]} & \textbf{Bit[0]} \\\hline
    保留 & lost\_source & format & anchored(0) \\\hline
\end{longtable}

从包格式可以看出，DMA\_0 和 DMA LOST 数据包均为 32 位。以下为各个域的含义：

\begin{itemize}
    \item \textbf{format} 表示数据包的类型：
    \begin{itemize}
        \item 0：DMA\_0 数据包
        \item 1：DMA LOST 数据包
    \end{itemize}
    \item \textbf{addr\_offset} 记录此次写操作的地址偏移，实际地址 = addr\_offset + \{\hyperref[regdesc:MEMMONITORLOGMINREG]{MEM\_MONITOR\_LOG\_MIN\_REG}[31:2]，2'b0\}。
    \item \textbf{dma\_source} 记录 DMA 访问来源，详见表 \ref{tab:assistdbg-dma-source}。
    \item \textbf{lost\_source} 记录 DMA LOST 数据包生成时舍弃的数据包：
    \begin{itemize}
        \item 0：未舍弃 DMA\_0 数据包
        \item 1：舍弃了 DMA\_0 数据包
    \end{itemize}
    \item \textbf{anchored} 表示该 32 位数据在数据包中的位置：
    \begin{itemize}
        \item 0：该 32 位数据是数据包的低 32 位
        \item 1：该 32 位数据是数据包的高 32 位
    \end{itemize}
\end{itemize}

模块内部 buffer 的数据宽度为 32 位。若使能 DMA\_0 通道组总线访问记录，并在某一时刻生成记录数据，DMA\_0 的数据包会缓存到 buffer 中。辅助调试模块会自动将缓存数据从 buffer 取出，并将取出的数据以 32 位数据宽度存入指定范围内存。

数据包解析流程如下：

\begin{itemize}
    \item 读取 \hyperref[fielddesc:MEMMONITORLOGMEMFULLFLAG]{MEM\_MONITOR\_LOG\_MEM\_FULL\_FLAG} 确定数据是否溢出配置范围。
    \begin{itemize}
        \item 如果未溢出，则数据读取的地址范围为 \hyperref[regdesc:MEMMONITORLOGMEMSTARTREG]{MEM\_MONITOR\_LOG\_MEM\_START\_REG}  \verb+~+ \hyperref[regdesc:MEMMONITORLOGMEMCURRENTADDRREG]{MEM\_MONITOR\_LOG\_MEM\_CURRENT\_ADDR\_REG} $-$ 4；
        \item 若溢出并且开启了 loop 模式，则数据读取的地址范围为 \hyperref[regdesc:MEMMONITORLOGMEMCURRENTADDRREG]{MEM\_MONITOR\_LOG\_MEM\_CURRENT\_ADDR\_REG} \verb+~+ \hyperref[regdesc:MEMMONITORLOGMEMENDREG]{MEM\_MONITOR\_LOG\_MEM\_END\_REG} 和 \hyperref[regdesc:MEMMONITORLOGMEMSTARTREG]{MEM\_MONITOR\_LOG\_MEM\_START\_REG} \verb+~+ \hyperref[regdesc:MEMMONITORLOGMEMCURRENTADDRREG]{MEM\_MONITOR\_LOG\_MEM\_CURRENT\_ADDR\_REG} $-$ 4；
        \item 若溢出并且未开启 loop 模式，则数据读取的地址范围为 \hyperref[regdesc:MEMMONITORLOGMEMSTARTREG]{MEM\_MONITOR\_LOG\_MEM\_START\_REG} \verb+~+ \hyperref[regdesc:MEMMONITORLOGMEMENDREG]{MEM\_MONITOR\_LOG\_MEM\_END\_REG}。
    \end{itemize}
    \item 从起始地址处开始读取数据并解析，每次读取 32 位。
\end{itemize}

数据解析完成之后需要通过置位 \hyperref[fielddesc:MEMMONITORCLRLOGMEMFULLFLAG]{MEM\_MONITOR\_CLR\_LOG\_MEM\_FULL\_FLAG} 清除 \hyperref[fielddesc:MEMMONITORLOGMEMFULLFLAG]{MEM\_MONITOR\_LOG\_MEM\_FULL\_FLAG} 标志位。


%%%%%%%%%%%%%%%%%%%%%%%%
%%% Register Summary %%%
%%%%%%%%%%%%%%%%%%%%%%%%

\newpage
\begin{landscape}
\hypertarget{dbgassist-reg-summ}{}
\section{寄存器列表}
\label{sec:dbgassist-reg-summ}

本小节中\textbf{总线记录配置寄存器}的地址为相对于存储器访问监控 1 (\textbf{TCM\_MEM\_MONITOR}) 基地址和存储器访问监控 2 (\textbf{PSRAM\_MEM\_MONITOR}) 基地址的地址偏移量（相对地址），详见小节 \ref{sec:dbgassist-mem-monitor-reg-summ}，其它寄存器的所有地址均为相对于总线访问监控(\textbf{BUS\_MONITOR}) 基地址的地址偏移量（相对地址），详见小节 \ref{sec:dbgassist-bus-monitor-reg-summ}，具体基地址请见章节 \ref{mod:sysmem} \textit{\nameref{mod:sysmem}} 中的表 \ref{tab:sysmem-base-address}。

请查看章节 \hyperref[glossary-access-types]{\textit{寄存器的访问类型}}，了解“访问”列缩写的含义。

\subsection{总线记录配置寄存器列表}
\label{sec:dbgassist-mem-monitor-reg-summ}

\subfile{\modulefiles/19-DBGASSIST-mem-monitor-reg-summ__CN}

\subsection{其他寄存器列表}
\label{sec:dbgassist-bus-monitor-reg-summ}

\subfile{\modulefiles/19-DBGASSIST-bus-monitor-reg-summ__CN}
\end{landscape}

%%%%%%%%%%%%%%%%%
%%% Registers %%%
%%%%%%%%%%%%%%%%%

\newpage
\section{寄存器}

本小节中\textbf{总线记录配置寄存器}的地址为相对于存储器访问监控 1 (\textbf{TCM\_MEM\_MONITOR}) 基地址和存储器访问监控 2 (\textbf{PSRAM\_MEM\_MONITOR}) 基地址的地址偏移量（相对地址），详见小节 \ref{sec:dbgassist-mem-monitor-reg}，其它寄存器的所有地址均为相对于总线访问监控(\textbf{BUS\_MONITOR}) 基地址的地址偏移量（相对地址），详见小节 \ref{sec:dbgassist-bus-monitor-reg}，具体基地址请见章节 \ref{mod:sysmem} \textit{\nameref{mod:sysmem}} 中的表 \ref{tab:sysmem-base-address}。

关于保留 (reserved) 域的处理，请查看章节 \hyperref[programming-reserved-field]{\textit{如何配置寄存器的保留域}}。

\subsection{总线记录配置寄存器}
\label{sec:dbgassist-mem-monitor-reg}

\subfile{\modulefiles/19-DBGASSIST-mem-monitor-reg__CN}

\subsection{其他寄存器}
\label{sec:dbgassist-bus-monitor-reg}

\subfile{\modulefiles/19-DBGASSIST-bus-monitor-reg__CN}

\end{document}
